\subsection{Executing declared event rules through a company program}
\label{sec:declarative-event-program}

A collection that works in isolation can still fail in a company program.
An earlier sale changes the claim population and its state schema; a report
advances the observation date; a year-end carry establishes the next opening
book. The method for a later event must be constructed from that current state.
Replaying a previously finished method cannot establish that these dependencies
are handled correctly.

\code{DeclarativeEventCompanyProgram/1} composes the explicit mechanisms.
The program supplies its opening book, typed memory definition, report mapping,
authored state invariant and one eligibility/construction pair per operation.
Event commands contain the full intent. An assessment also supplies an ordered
list covering its current claim population. The program contains neither native
event callbacks nor precomputed event results. Separate hashes bind the selected
rules and invariant; program identity also retains the order of event inputs.

Before each command, the runner checks identity uniqueness, the typed book and
the selected state invariant. For an event, ordered eligibility runs first.
Only an eligible event reaches the parameter or ordered-row constructor, which
uses the current book and full intent to produce the complete transition
request. The unchanged typed transition then computes postings, flows and
memory effects. A report uses the current period's opening and the book observed
at its stated cutoff. A carry uses the existing closing and carry evaluator,
with the explicitly selected policy preserving every memory leaf.

The book advances only after the resulting state also passes its type and
invariant checks. Consider an authored sale whose postings increase receivables
but whose inserted claim incorrectly records zero remaining amount. The typed
transition can propose that state because each value and destination has a
valid type. The separately selected claim-total invariant rejects it: claim
amounts no longer equal the receivables balance. The runner retains the prior
current and opening books, records the refused command, and leaves the complete
remaining suffix unexecuted. It adds no compensating journal or financing event.
This is in-memory refusal behaviour, not a claim of durable database rollback.

The two retained eleven-command programs preserve all complete native command
requests and results, including reports and carry archives, along with their
before/after state identities and final books. Two further eleven-command
programs start from the same declared openings but create three new claims each,
using changed keys, customers and amounts. Their report, year-end carry and
following-year assessment also agree. All native lifecycle functions are
disabled during candidate execution. The saved result is checked again after
serialization; no native after-state is supplied to the candidate.

Twenty-six adverse programs include the original eighteen refusal cases,
combined unknown/approval/missing-claim conditions, and negative or excessive
collections that fail financial preconditions after eligibility has passed.
The latter retain the complete refused request and transition. Additional
controls check stale selections, changed input order, duplicate identities,
incomplete assessment order, a changed carry policy and the invariant-invalid
sale. Each stopped run preserves its successful prefix, current opening,
reports and unexecuted suffix.

The comparison concerns complete substantive command payloads. The candidate
adds explicit eligibility, construction and invariant records; the native
result carries different implementation and observation records. Those
differences are retained and identified, rather than removed and described as
whole-result equality. Native observations are separately recomputed from the
candidate states. Shared transition, report and carry evaluators remain the
owners of their respective financial calculations.

Opening admission remains a separate native prerequisite for this composer
comparison. The original validation checks the retained opening/program before
native functions are disabled. The next extension declares those opening
conditions and binds them to the exact program selected for execution. The
earlier company-source frontend retains its own native-contract target.
General opening populations, a complete financial language, accounting validity
and independent implementation acceptance remain separate questions.

\begingroup\raggedright
The four complete programs, adverse cases and full comparisons are retained
under \path{data/accounting-master/declarative-event-program-20261026/}.
\par
\endgroup
